\documentclass[a4paper,11pt]{article}
\usepackage[utf8]{inputenc}
\usepackage[T1]{fontenc}
\usepackage{lmodern}
\usepackage[margin=2.5cm]{geometry}
\usepackage{graphicx}
\usepackage{booktabs}
\usepackage{array}
\usepackage{float}
\usepackage{hyperref}
\usepackage{enumitem}

\title{Segmentação e Reconhecimento de Rótulos em Embalagens de Frango\\
\large Documento técnico: arquitetura, requisitos, dados, testes e resultados}
\author{Bruno Silveira (@cenourin)}
\date{}

\begin{document}
\maketitle
\tableofcontents

\section{Visão geral}

Este documento descreve a arquitetura e os resultados do projeto de Processamento Digital de
Imagens (PDI) composto por duas etapas sobre o mesmo conjunto de fotos de embalagens de produtos
avícolas, capturadas em uma esteira industrial com câmera fixa:

\begin{itemize}[nosep]
  \item \textbf{T1 --- Segmentação}: localizar e recortar a região do rótulo com o nome do
        produto em cada imagem de embalagem.
  \item \textbf{T2 --- Reconhecimento}: identificar a qual classe de produto pertence um
        segmento gerado no T1, comparando-o com um template de referência por classe.
\end{itemize}

Ambas as etapas usam exclusivamente técnicas clássicas de PDI (espaços de cor, limiarização,
morfologia matemática, filtragem espacial, transformações geométricas, descritores locais de
características) --- sem uso de aprendizado de máquina ou classificadores treinados em nenhuma
das duas etapas.

\section{Requisitos}

\subsection{Requisitos funcionais}

\begin{enumerate}[nosep]
  \item RF01 --- Dado um caminho de imagem de embalagem, o pipeline de T1 deve retornar um
        recorte quadrado contendo a região do nome do produto, ou nenhum recorte quando não há
        candidato confiável.
  \item RF02 --- O pipeline de T1 deve tratar dois layouts físicos distintos de embalagem
        (bandeja e selado), cada um com um pipeline próprio.
  \item RF03 --- Dado um segmento de T1 (ou um recorte manualmente corrigido) e um conjunto de
        templates gold (um por classe), o pipeline de T2 deve retornar a classe de maior
        correspondência ou \emph{unknown} quando nenhum template atinge o limiar mínimo de
        confiança.
  \item RF04 --- Ambos os pipelines devem processar o dataset inteiro em lote e produzir um
        resumo por classe (sucesso/falha em T1; acerto/não identificado/falso positivo em T2).
\end{enumerate}

\subsection{Requisitos não funcionais}

\begin{enumerate}[nosep]
  \item RNF01 --- \textbf{Zero alteração de código em tempo de avaliação}: os parâmetros dos
        pipelines são centralizados (dicionário \texttt{CFG} em T1; \texttt{RecognitionConfig}
        em T2) e devem generalizar para um conjunto de imagens novo, não apenas para as imagens
        vistas durante o desenvolvimento.
  \item RNF02 --- \textbf{Nenhuma biblioteca de segmentação ou classificação baseada em IA/ML}:
        apenas operações clássicas de PDI (OpenCV: limiarização, morfologia, contornos,
        SIFT/ORB + FLANN para correspondência de descritores).
  \item RNF03 --- \textbf{Falso positivo é mais custoso que falso negativo}: em ambas as etapas,
        a calibração de parâmetros prioriza reduzir a taxa de falso positivo mesmo à custa de
        recall (mais recortes/classificações rejeitados como inválidos ou \emph{unknown}).
\end{enumerate}

\section{Dados}

O dataset (\texttt{Train\_and\_Validation/}) contém 18 classes de produto, cada uma em uma pasta
nomeada \texttt{<código>\_<Nome\_Do\_Produto>}, com aproximadamente 50 fotos por classe tiradas
por um rig de câmera fixo sobre a esteira de produção. As classes se dividem em dois tipos de
embalagem física, que exigem pipelines de segmentação diferentes:

\begin{itemize}[nosep]
  \item \textbf{Bandeja}: rótulo como região destacada e não fixa dentro do quadro (7 classes
        usadas na etapa de reconhecimento do T2).
  \item \textbf{Selado}: rótulo em uma faixa horizontal fixa da embalagem, com uma região de
        FPS/timestamp sobreposta que precisa ser mascarada (11 classes).
  \end{itemize}

O dataset completo não é versionado neste repositório (ver \texttt{README.md}) por ser grande
demais e por conter imagens de produto de uma marca comercial (Super Frango) cujos direitos não
pertencem ao autor deste trabalho.

\section{Arquitetura --- T1 (Segmentação)}

Pipeline linear por imagem, cada etapa uma função pura operando sobre a saída da etapa anterior,
com todos os parâmetros centralizados em um dicionário de configuração em vez de espalhados como
literais no código.

\subsection{Pipeline bandeja (\texttt{segmentador\_colab.ipynb})}

\begin{enumerate}[nosep]
  \item \texttt{extrair\_brilho} --- conversão BGR→LAB e CLAHE no canal L, para compensar
        iluminação industrial desigual.
  \item \texttt{suavizar} --- filtro de mediana (remoção de ruído).
  \item \texttt{binarizar} --- limiarização de Otsu com deslocamento manual
        (\texttt{thresh\_offset}): Otsu sozinho tende a escolher um limiar entre o texto escuro
        do rótulo e o fundo mais claro, então é ajustado empiricamente.
  \item \texttt{aplicar\_roi} --- zera uma margem fixa de borda (esteira/paredes laterais ficam
        fora de \texttt{roi\_left/right/top/bot}).
  \item \texttt{morfologia} --- fechamento (junta o texto do rótulo em um blob) seguido de
        abertura (remove blobs de ruído pequenos).
  \item \texttt{classificar\_regioes} --- extração de contornos com filtros geométricos (área,
        proporção largura/altura, altura mínima, solidez), anotando o motivo de rejeição de
        cada região descartada.
  \item \texttt{nms} --- supressão de sobreposição por IoU entre caixas candidatas; apenas a
        maior caixa aceita é mantida.
  \item \texttt{square\_crop} --- recorte quadrado centralizado na caixa vencedora (lado =
        max(largura, altura), sem redimensionamento).
\end{enumerate}

\subsection{Pipeline selado (\texttt{segmentador\_selado.ipynb})}

Layout físico diferente do bandeja (rótulo em faixa horizontal fixa em vez de região solta):
alongamento de contraste por percentil → recorte de ROI com máscara fixa da região de
FPS/timestamp → morfologia top-hat → limiarização adaptativa/fixa → componentes conectados
filtrados por área → \emph{clustering} espacial de componentes para reagrupar texto
multi-letra em uma única região, preferindo o agrupamento de proporção larga (horizontal), que é
a orientação do texto nesse tipo de embalagem.

\section{Arquitetura --- T2 (Reconhecimento)}

\begin{enumerate}[nosep]
  \item \textbf{Templates}: um recorte manual por classe (região do rótulo com maior potencial
        discriminativo), armazenado em \texttt{t2/templates-gold/}.
  \item \textbf{Extração de características}: SIFT (\texttt{cv2.SIFT\_create}), com ORB
        disponível como alternativa via configuração.
  \item \textbf{Casamento}: \texttt{cv2.FlannBasedMatcher}, com teste de razão de Lowe aplicado
        nos dois sentidos (cross-check).
  \item \textbf{Verificação geométrica}: RANSAC com transformação afim sobre os matches válidos
        --- descarta correspondências espúrias em elementos repetidos entre rótulos (marca,
        palavra "CONGELADO", tabela de peso), a principal fonte de falso positivo nesse tipo de
        matching.
  \item \textbf{Score}: fração de matches válidos (inliers) sobre o número de pontos de
        interesse do template. O template de maior score vence, desde que o score atinja o
        limiar \texttt{min\_match\_frac}; caso contrário, o segmento é reportado como
        \emph{unknown}.
\end{enumerate}

Detalhamento de parâmetros e a varredura de calibração de \texttt{min\_match\_frac} estão em
\texttt{docs/relatorio-t2.tex}.

\section{Metodologia de teste}

\subsection{T1}

Execução em lote sobre o dataset completo, com contagem por classe de sucesso (recorte gerado) e
falha, e motivo de rejeição registrado por região descartada (área insuficiente, proporção fora
da faixa, altura mínima não atingida, baixa solidez). Falso positivo (recorte sem o texto do
produto, ou com o texto errado) é tratado como erro mais grave que detecção perdida, conforme o
enunciado da disciplina.

\subsection{T2}

Avaliação sobre 350 imagens (7 classes bandeja × 50 imagens), usando os segmentos gerados no T1
e corrigidos manualmente quando o T1 não gerou recorte. Métricas: acurácia (classificação
correta), taxa de não identificado (\emph{unknown}, quando nenhum template atinge o limiar) e
taxa de falso positivo (classificado com produto errado). O parâmetro \texttt{min\_match\_frac}
foi calibrado por varredura, buscando o menor valor que mantém a taxa de falso positivo baixa
sem descartar recall demais --- ver tabela de calibração em
\texttt{docs/relatorio-t2.tex}.

\section{Análise de resultados}

\subsection{T1}

Amostras de resultado (uma imagem por classe, grids com os recortes gerados) estão em
\texttt{docs/resultados-exemplo/bandeja/} (7 classes) e
\texttt{docs/resultados-exemplo/selado/} (11 classes) --- as 18 classes do dataset. A saída
completa do pipeline (todas as imagens, todas as classes) não é versionada; pode ser
regenerada rodando os notebooks sobre o dataset local.

\subsection{T2}

Com \texttt{min\_match\_frac = 0,08}: acurácia global de 83,4\% (292/350), 48 segmentos não
identificados e 10 falsos positivos (2,9\%) sobre as 350 imagens avaliadas. O desempenho varia
por classe --- produtos com texto curto e exclusivo (Coração, Filé de Peito) atingem 100\% de
acerto, enquanto classes com rótulos de menor contraste ou mais elementos repetidos entre
embalagens (Filé de Coxas e Sobrecoxas com Pele, Coxinhas das Asas) têm mais segmentos
descartados como \emph{unknown}. Detalhamento completo por classe e a matriz de confusão estão
em \texttt{docs/relatorio-t2.tex} e \texttt{docs/matriz\_confusao\_t2.png}.

\section{Aplicação em um caso real}

O pipeline desenvolvido é um protótipo direto de um sistema de \textbf{controle de qualidade em
linha de produção} para indústria alimentícia: uma câmera fixa sobre a esteira de embalagem
captura cada unidade, o T1 localiza e recorta o rótulo, e o T2 confirma se o rótulo corresponde
ao produto esperado para aquela linha/turno. Um \emph{unknown} ou uma classificação divergente
da SKU configurada pode disparar um alerta ou desviar a unidade da linha antes do
despacho, evitando o erro de expedir produto com rótulo trocado.

A prioridade de projeto adotada em ambas as etapas --- minimizar falso positivo mesmo à custa de
recall --- reflete diretamente o custo assimétrico desse cenário real: deixar passar uma
embalagem sem verificação automática (falso negativo, cai para inspeção manual) é muito mais
barato que confirmar automaticamente um rótulo errado como correto (falso positivo, produto
sai da fábrica mal rotulado).

Limitações para uso em produção: os parâmetros de segmentação foram calibrados para as condições
de iluminação e o rig de câmera específicos deste dataset; uma implantação real exigiria
recalibração para cada linha/câmera, e o tempo de processamento por imagem (SIFT + RANSAC) precisaria
ser validado contra a cadência real da esteira.

\section{Direitos autorais}

As imagens do dataset utilizado neste trabalho retratam embalagens e produtos da marca
\textbf{Super Frango}, cujos direitos pertencem aos seus respectivos titulares. As imagens foram
disponibilizadas exclusivamente para fins acadêmicos nesta disciplina. Este documento e o
repositório associado versionam apenas o código de processamento de imagem desenvolvido pelo
autor; nenhuma reivindicação de direito é feita sobre a marca, o produto ou as imagens
originais do dataset, e o autor não se responsabiliza pelo uso de tais imagens fora do escopo
acadêmico deste trabalho.

\end{document}
